\subsection{Sheafification}

There is a natural map
$$ \F(U) \to \F_p , $$
for every $p\in U$ given by mapping $s\in\F(U)$ to $(U,s)\in\F_p$.
Together these assimilate into a map
$$ \F(U) \to \prod_{p\in U}\F_p . $$

\bprop[name={inctostalk}]

    The natural map above is injective when $\F$ is a sheaf.

\eprop

\bproof

    Suppose that $s,t\in\F(U)$ such that $(U,s)=(U,t)$ in $\F_p$.
    Then there exists a neighborhood $p\in W_p\subseteq U$ such that $s|_{W_p}=t|_{W_p}$.
    If this is true for all $p$, then $W_p$ form an open cover of $U$ and by the identity axiom, $s=t$.
    \qed

\eproof

\bdefn

    An element $(s_p)_p\in\prod_{p\in U}\F_p$ consists of {\emphcolor compatible germs} if for all $p\in U$, $s_p$ has a representative
    $\tilde s_p\in\F(U_p)$ for $p\in U_p\subseteq U$ such that the germ of $\tilde s_p$ at $q\in U_p$ is $s_q$ for all $q\in U_p$.

\edefn

Equivalently, this says that there is an open cover $\set{U_i}$ of $U$ and sections $f_i\in\F(U_i)$ such that if $p\in U_i$ then $s_p$ is the germ
of $f_i$ at $p$.
If $s$ is a section of $\F$ over $U$, then $(s)_p$ consists of compatible germs.

\bprop[name={inctostalkim}]

    The natural map $\F(U)\to\prod_{p\in U}\F_p$'s image is precisely the compatible germs.

\eprop

\bproof

    We already showed that every tuple in its image consists of compatible germs, now we show the converse.
    Suppose $(s_p)_p$ consists of compatible germs, so let $\tilde s_p\in\F(U_p)$ be a representation as in the definition.
    Then $U_p$ cover $U$, and $\tilde s_p|_{U_p\cap U_q}$'s germ at $q$ is $s_q$ so that $\tilde s_p,\tilde s_q$ agree on a neighborhood of $q$
    and by symmetry in a neighborhood of $p$ and $q$.
    \qed

\eproof

\bprop[name={stalksdeterminemorphism}]

    Suppose $\F$ is a presheaf of sets and $\G$ a sheaf of sets, and $\phi_1,\phi_2\colon\F\to\G$ are two morphisms.
    Further suppose that $\phi_1,\phi_2$ induce the same map on each stalk; show that $\phi_1=\phi_2$.

\eprop

\bproof

    Consider the following diagram for $U\subseteq X$ open,

    \commsq{\F(U)&\G(U)\cr
        \prod_{p\in U}\F_p&\prod_{p\in U}\G_p}
        {\phi_i}{\prod\phi_{i,p}}{}{\iota}

    for $i=1,2$.
    It is easily seen that this diagram commutes, and because $\phi_{1,p}=\phi_{2,p}$ for all $p\in U$ it follows that
    $$ \iota\circ\phi_1 = \iota\circ\phi_2 $$
    and since $\iota$ is injective by \refmath[proposition]{inctostalk}, we get that $\phi_1=\phi_2$ as desired.
    \qed

\eproof

\bprop[name={isoonstalks}]

    A morphism of sheaves of sets is injective iff it induces an injection on all stalks.
    Similarly it is an isomorphism iff it induces an isomorphism on all stalks.

\eprop

\bproof

    If $\phi\colon\F\to\G$ is an isomorphism, then by functorality $\phi_p\colon\F_p\to\G_p$ is as well for all $p$.
    Similarly if it is an injection then $\phi_p(s_p)=\phi_p(t_p)$ then $\phi(s)$ and $\phi(t)$ are equal in a subset $V\subseteq U$ so that
    $\phi(s|_V)=\phi(t|_V)$ and by injectivity this means that $s|_V=t|_V$ and thus $s_p=t_p$.

    Conversely consider the commutative diagram

    \diagram{
        \F(U)&\G(U)\cr
        \prod_{p\in U}\F_p&\prod_{p\in U}\G_p
    }{
        \diagarrow{from={1,1}, to={1,2}, text=\phi, y distance=.25cm}
        \diagarrow{from={2,1}, to={2,2}, text=\sim, y distance=.25cm}
        \diagarrow{from={1,1}, to={2,1}, left cap={(}, text={\iota_\F}, x distance=-.25cm}
        \diagarrow{from={1,2}, to={2,2}, left cap={(}, text={\iota_\G}, x distance=.25cm}
    }

    where the bottom horizontal arrow is the injection/isomorphism $\prod_{p\in U}\phi_p$.
    Then $\iota_\G\circ\phi$ is injective, and so $\phi$ is.

    Now to show that $\phi$ is surjective, let $s\in\G(U)$.
    Then $\iota_\G(s)\in\prod_U\G_p$ consists of compatible germs, and mapping it back to $\prod_U\F_p$ we see that it still
    consists of compatible germs.
    Thus it lies in the image of $\iota_\F$ by \refmath[proposition]{inctostalkim}, and thus has a preimage in $\F(U)$.
    Because all of the maps are injective, this is the preimage of $s$ in $\F(U)$ as desired.
    \qed

\eproof

\bnote

    It is not true that sheaves with isomorphic stalks are isomorphic: there may not be an isomorphism between the sheaves which induces these
    isomorphisms.

\enote

\bthrm[title={Sheafification}, name={sheafification}]

    The inclusion $\Set_X\inj\Set_X^\pre$ has a left adjoint, called {\emphcolor sheafification}: $\sh\colon\Set_X^\pre\to\Set_X$.

\ethrm

\bproof

    We need to show for any presheaf $\F$ there exists a sheaf $\F^\sh$ and morphism $\sh\colon\F\to\F^\sh$ such that for every sheaf $\G$ and
    morphism $g\colon\F\to\G$ there exists a unique morphism $f\colon\F^\sh\to\G$ such that the below diagram commutes:

    \diagram{
        \F&\F^\sh\cr
        &\G
    }{
        \diagarrow{from={1,1}, to={1,2}, text=\sh, y distance=.25cm}
        \diagarrow{from={1,1}, to={2,2}, text=g, x distance=-.25cm}
        \diagarrow{from={1,2}, to={2,2}, text={\exists!f}, x distance=.5cm}
    }

    We construct $\F^\sh$ explicitly.
    For $U\subseteq X$ open, define
    $$ \F^\sh(U) = \set{(f_p\in\F_p)_{p\in U}}[\mltext{for all $p\in U$ there exists a neighborhood $p\in V\subseteq U$ and $s\in\F(V)$\cr
    such that $s_q=f_q$ for all $q\in V$}] , $$
    where by $s_q$ we mean the germ of $s$ in $\F_q$.
    For $U\subseteq V$ there is an obvious map $\F^\sh(V)\to\F^\sh(U)$ mapping $(f_p)_{p\in V}\mapsto(f_p)_{p\in U}$; it is easily seen that
    $\F^\sh$ is a presheaf.

    Note that the construction of $\F^\sh(U)$ is just the set of all families of compatible germs over $U$, which we discussed in the
    previous few results.

    Now let $U\subseteq X$ be open and $\set{U_i}_i$ an open cover of $U$.
    Then if we have $(f^i_p)_{p\in U_i}\in\F^\sh(U_i)$ for all $i$ which are compatible, this means $f^i_p=f^j_p$ for all $p\in U_i\cap U_j$.
    Then we must define $(f_p)_{p\in U}$ by $f_p=f^i_p$ for any $i\in U_i$ for which $p$ is a member of.
    This is well-defined precisely by compatibility, and it is a member of $\F^\sh(U)$ because for $p\in U$ there is a $p\in U_i\subseteq U$
    and then there is a $p\in V\subseteq U_i$ and $s\in\F(V)$ such that $s_q=f^i_q=f_q$ for all $q\in V$ as desired.
    Thus $\F^\sh$ is a sheaf.

    Now we must demonstrate that $\F^\sh$ satisfies the universal property of a left adjoint.
    There is an obvious map $\sh\colon\F\to\F^\sh$ which maps $s\in\F(U)$ to $(s_p\in\F_p)_{p\in U}\in\F^\sh(U)$ (this is a section because
    $s$ satisfies the condition for it to be in $\F^\sh(U)$).
    Let $g\colon\F\to\G$ be a morphism, and for $U\subseteq X$ consider the following diagram

    \diagram{
        \F(U)&\F^\sh(U)&\prod_{p\in U}\F_p\cr
        &\G(U)&\prod_{p\in U}\G_p
    }{
        \diagarrow{from={1,1}, to={1,2}, text=\sh, y distance=.25cm}
        \diagarrow{from={1,1}, to={2,2}, text=g, x distance=-.25cm}
        \diagarrow{from={1,2}, to={1,3}, text={\iota_\F}, y distance=.25cm, left cap={(}}
        \diagarrow{from={2,2}, to={2,3}, text={\iota_\G}, y distance=-.25cm, left cap={(}}
        \diagarrow{from={1,3}, to={2,3}, text={\prod_pg_p}, x distance=.25cm}
    }

    Note that the image of $\prod_pg_p\circ\iota_\F$ consists of compatible germs in $\prod_{p\in U}\G_p$ and so we can fill this diagram
    in with a unique arrow $f\colon\F^\sh(U)\to\G(U)$.
    It is easy to see by uniqueness that these $f$ assemble into a morphism $\F^\sh\to\G$.
    \qed

\eproof

\bnote

    In the language of category theory, this theorem says that $\Set_X\subseteq\Set_X^\pre$ is a {\emphcolor reflective subcategory}.
    Note that this proof works for other concrete categories like $\Ab,\Grp,\Ring$, etc.
    It can be shown that for sufficiently nice $\C$ (e.g.\ complete and cocomplete), sheafification exists for $\C_X\subseteq\C_X^\pre$.

    This justifies the fact that inclusion $\C_X\subseteq\C_X^\pre$ preserves limits: it is a right adjoint.
    This explains why the presheaf kernel of a sheaf morphism is a sheaf kernel for example.

\enote

\bprop

    The canonical morphism $\F\to\F^\sh$ induces an isomorphism on stalks.

\eprop

\bproof

    This is easily seen by an explicit description of $\F^\sh$'s stalks.
    \qed

\eproof

\bremark[title={The space of sections of a presheaf}]

    Suppose $\F$ is a presheaf of sets on $X$.
    We will construct a topological space $F$ along with a continuous map $\pi\colon F\to X$ as follows.
    As a set define $F$ to be the disjoint unions of $\F$'s stalks:
    $$ F = \coprod_{p\in X}\F_p . $$
    Now for $s\in\F(U)$ let $F_s=\set{(p,s_p)}[p\in U]\subseteq F$; then the collection of all $F_s$ where $s\in\F(U)$ for arbitrary $U\subseteq X$
    forms a base for a topology on $F$.
    Indeed, for $s\in\F(U)$ and $t\in\F(V)$ we see that $F_s\cap F_t$ is the union of all $F_r$ where $r=s|_W=t_W$ for all $W\subseteq U\cap V$
    where $s|_W=t|_W$.

    The map $\pi\colon F\to X$ is the obvious projection $(p,s_p)\mapsto p$.
    For $U\subseteq X$ open we have that $\pi^{-1}U$ is the union of all $F_s$ where $s\in\F(U)$ so that $\pi$ is continuous.
    Furthermore $\pi$ is an open map: for $s\in\F(U)$ it maps $F_s$ to $U$.
    Moreso $\pi|_{F_s}$ is a bijection, so that $\pi$ is a local homeomorphism.

    The space $F$ is called the {\emphcolor space of sections}, or {\emphcolor \'etal\'e space} of $\F$.

    We can give an alternative construction for the sheafification of $\F$ as the sheaf of sections of $\pi\colon F\to X$
    (that is for $U\subseteq X$ we define $\F^\sh(U)$ to be the set of continuous maps $\sigma\colon U\to F$ where $\pi\circ\sigma=\id_U$).
    Let us denote by $\F^+$ this new construction.

    For $U\subseteq X$ open we define $\phi\colon\F^\sh(U)\to\F^+(U)$ by mapping $(f_p)_{p\in U}$ to $\sigma\colon U\to F$, $p\mapsto f_p$.
    Its inverse is given by mapping $\sigma$ to $(\sigma(p))_p$.

\eremark

\bexam

    Let $S$ be a set and consider the constant presheaf $\ul S_\pre$.
    Then its sheafification is the constant {\it sheaf\/}: $(\ul S_\pre)^\sh=\ul S$.

    We can show this with the universal property: $\sh\colon\ul S_\pre\to\ul S$ is given by inclusion (constant maps are locally constant).
    Then for a morphism $g\colon\ul S_\pre\to\F$ where $\F$ is a sheaf, we define $\bar g\colon\ul S\to\F$ by mapping a locally constant
    $f\colon U\to S$ to the gluing of $(\bar g(f|_{U_i}))_i$ (where $U_i\subseteq U$ are its connected components so that $f|_{U_i}$ is a constant).

\eexam

Let $\phi\colon\F\to\G$ be a morphism of sheaves, then recall the following equivalences (see \refmath[remark]{sheafinclim} and
\refmath[proposition]{isoonstalks}):
\benum
    \item $\phi$ is monic as a morphism of sheaves,
    \item $\phi$ is monic as a morphism of presheaves,
    \item $\phi(U)$ is monic for all $U\subseteq X$,
    \item $\phi_p$ is monic for all $p\in X$.
\eenum

A similar equivalence does not hold true for epimorphisms.
This boils down to the fact that inclusion $\C_X\subseteq\C_X^\pre$ does not necessarily preserve colimits, and thus not epimorphisms.
So while a presheaf morphism is epic iff its components are, the same does not hold for sheaf morphisms.
But it is true that epimorphisms are determined by stalks.

\bprop[name={epionstalks}]

    Let $\phi\colon\F\to\G$ be a morphism of sheaves.
    Then it is an epimorphism (in the category of sheaves) iff $\phi_p\colon\F_p\to\G_p$ is an epimorphism for all $p\in X$.

\eprop

\bproof

    If $\phi_p$ is epic for all $p$, then suppose we have two sheaf morphisms $\psi_1,\psi_2\colon\G\to\c H$ such that
    $\psi_1\circ\phi=\psi_2\circ\phi$.
    Then taking stalks we see that for all $p$, $(\psi_1)_p\circ\phi_p=(\psi_2)_p\circ\phi_p$.
    Since $\phi_p$ is epic this means $(\psi_1)_p=(\psi_2)_p$ for all $p$ and by \refmath[proposition]{stalksdeterminemorphism} this means
    $\psi_1=\psi_2$ so that $\phi$ is epic.

    Conversely let $p\in X$ and suppose there are maps $f,g\colon\G_p\to X$ such that $f\circ\phi_p=g\circ\phi_p$.
    We can define maps $\bar f,\bar g\colon\G\to p_*X$ by $\bar f_U(s)=f(s_p)$ for $p\in U$ and $\bar f_U=*$ otherwise.
    It is easy to see that $\bar f\circ\phi=\bar g\circ\phi$ and so $\bar f=\bar g$ which can easily be shown to mean $f=g$.
    \qed

\eproof

\bnote

    In \refmath[remark]{stalkskyscraperadj} will show that there is an adjunction between taking stalks and the skyscraper
    sheaf: $\bul_p\dashv p_*\bul$.
    This was used in the above proof, and it is easy to see then how to construct this.

\enote

\bexam[name={expsequence}]

    We will show that an epimorphism of sheaves need not be epic on each of its components.

    Let $\O_\bC$ be the sheaf of holomorphic functions and $\O_\bC^*$ be the sheaf of invertible (nowhere zero) holomorphic functions.
    These are sheaves of Abelian groups, and we then have maps of sheaves
    $$ 0 \to \ul\bZ \xtos{\times2\pi i} \O_\bC \xtos{\exp} \O_\bC^* \to 1 . $$
    We will soon show that this is an exact sequence of sheaves of Abelian groups.

    In particular $\exp\colon\O_\bC\to\O_\bC^*$ is an epimorphism as we can show directly.
    For any $z\in\C$ the morphism of stalks $\exp_z$ is surjective because every invertible holomorphic function locally has a logarithm.
    But $\exp$ is not an epimorphism on each of its components, for example $\id\in\O_\bC^*(\bC-0)$ has no global logarithm.

\eexam

